\section{Resumen}

Los mercados de criptoactivos combinan plataformas centralizadas (CEX) con creadores automáticos de mercado (AMM), cuyo precio depende de las reservas y de las operaciones ejecutadas. En los mercados reales, el arbitraje está condicionado por los costos, el capital y la liquidez; en este estudio su capacidad se representa con un límite especificado por tick, sin modelar comisiones ni restricciones de capital vinculantes. Este trabajo estudia cómo un AMM de producto constante sigue un precio CEX exógeno después de una perturbación, en un micromundo con agentes de ruido, tendencia, arbitraje y provisión de liquidez, y evalúa un mapa dinámico discreto de dos variables como descripción reducida. Se contrastan tres hipótesis con un protocolo confirmatorio fijado de antemano: semillas y celdas de confirmación disjuntas de las de exploración, inferencia pareada por semilla y corrección de Holm.

Con 300 semillas de confirmación y 13500 ejecuciones, la regla de arbitraje activo especificada reduce el área de exceso de error de 2.36 a 0.10 dentro del simulador ensayado (H1; diferencia media de −2.26; p de Holm 0.00015), y la capacidad gradúa la recuperación: 56.8 por ciento con capacidad baja y 100 por ciento con capacidad alta. El retiro adaptativo de liquidez aumenta en promedio la respuesta máxima en 0.011 (H2), pero solo con la menor profundidad; con profundidades mayores el efecto es nulo. El mapa reducido supera a la persistencia (H3), aunque no a un autorregresivo escalar, y su clasificación de regímenes fuera de muestra es débil. Un análisis suplementario de sensibilidad con 40\,500 trayectorias encontró que el 31.2 por ciento de las trayectorias adaptativas que se recuperaron salió después de la banda; la amplificación con profundidad baja desapareció al escalar las órdenes con las reservas y bajó de 0.0346 a 0.0147 unidades logarítmicas al desactivar el pánico. Estos resultados sintéticos y condicionados sugieren mediciones y preguntas candidatas para futuros estudios empíricos con reproducción histórica de eventos y pruebas de estrés; no validan procedimientos de monitoreo ni umbrales operativos. Un piloto independiente de disponibilidad de datos, de seis minutos, contiene tres intercambios públicos de un DEX; es solo descriptivo y no contrasta las hipótesis. Los hallazgos confirmatorios carecen de validación externa y no se alcanzó la precisión objetivo de H1.

\textbf{Palabras clave:} simulación basada en agentes; arbitraje; liquidez; seguimiento de precios; mercados descentralizados; sistemas dinámicos discretos
